\section{本章小结}
\label{sec:opt-smooth-summary}

章首的问题是：在$f(\bm x)=\tfrac12(x_1^2+16x_2^2)$上，只凭一个点的目标值和梯度，怎样决定方向与步长，怎样保证下降，何时值得读取曲率，返回的点有多准、花了多少代价。本章的回答围绕一条线索：在当前点用局部模型代替目标，并时刻知道模型与真实目标差多少。

梯度Lipschitz条件给出了最基本的模型：一个曲率为$L$、盖住目标的二次上界。它的最小点就是步长$1/L$的梯度步，下降引理把“沿负梯度走足够小的一步”变成可计算的保证。由此得到三层结论：$L$光滑且下有界时，梯度下降使历史最佳梯度以$O(T^{-1/2})$减小；再加凸性与最小点存在，最后一点的目标差以$O(1/T)$减小；满足PL条件时目标差线性收敛。回溯线搜索在不知道$L$时用实际函数值完成同样的工作。统一常数$L$的代价是对平坦方向过于保守：共同例中曲率相差$16$倍，固定步长只能在“之”字与缓慢之间折中。

改进有两条路。一条是改变坐标或记住历史：固定预条件等价于换坐标后的梯度下降，好坏由变换后的曲率决定；重球法在二次目标上把依赖$\kappa$的收缩变为依赖$\sqrt\kappa$，Nesterov方法在一般凸光滑目标上把$O(1/T)$变为$O(1/T^2)$，代价是目标值不再单调。另一条是把曲率放进模型：牛顿法在非退化最小点附近二次收敛，但需要Hessian正定、需要线搜索控制远处的步长，每轮代价高；修正牛顿、信赖域与立方正则化分别用改曲率、限半径、加立方项的方式保证模型可用；Newton--CG只用Hessian向量积求解模型，拟牛顿用梯度差拼出曲率，它们降低每轮代价，也相应改变了能保证的局部速度。

非凸目标上，梯度小只说明接近某个驻点。严格鞍点处有负曲率方向，可以用二阶信息直接下降，也可以用扰动梯度下降在有限步内高概率地离开；但近似二阶驻点仍不等于局部最小点，更不等于全局最小点。比较算法时，误差对象、返回点、条件与成本口径必须一致：局部收敛阶不代表全程成本，每轮便宜也不代表总成本低。计算误差给可达精度设下了下限，而数据误差引起的解变化由问题自身的曲率下界决定，再精确的求解也无法消除。

表~\ref{tab:opt-smooth-summary}汇总了本章主要方法每轮使用的信息、所需条件与代价。

\begin{table*}[htbp]
  \centering
  \small
  \begin{tabularx}{\linewidth}{@{}>{\raggedright\arraybackslash}p{0.15\linewidth}>{\raggedright\arraybackslash}p{0.2\linewidth}>{\raggedright\arraybackslash}X>{\raggedright\arraybackslash}p{0.2\linewidth}@{}}
    \toprule
    方法 & 每轮使用的信息 & 主要保证及其条件 & 每轮主要代价（稠密$d$维） \\
    \midrule
    梯度下降 & 梯度；回溯时加函数值 & $L$光滑：残差$O(T^{-1/2})$；凸：目标差$O(1/T)$；PL：线性 & 一次梯度，$O(d)$存储 \\
    固定预条件 & 梯度与固定的$\bm P$ & 等价于新坐标中的梯度下降，速度由变换后曲率决定 & 应用$\bm P$的代价 \\
    重球法 & 梯度与上一步位移 & 正定二次、谱界已知：因子$\tfrac{\sqrt\kappa-1}{\sqrt\kappa+1}$ & 一次梯度，多存一个向量 \\
    Nesterov方法 & 前瞻点的梯度 & 凸、$L$光滑、最小点存在：目标差$O(1/T^2)$ & 一次梯度，多存一个向量 \\
    牛顿法 & 梯度与Hessian & 非退化最小点邻域：二次收敛；远处需线搜索 & $O(d^3)$分解，$O(d^2)$存储 \\
    信赖域 & 梯度与模型矩阵 & Cauchy下降、模型有界、下有界：梯度趋零 & 子问题求解，拒绝时重解 \\
    Newton--CG & 梯度与Hessian向量积 & 局部阶由强迫序列决定 & $k$次Hessian向量积，$O(d)$存储 \\
    BFGS／L-BFGS & 梯度与历史位移、梯度差 & 曲率条件下保持正定；完整BFGS局部超线性 & $O(d^2)$／$O(rd)$ \\
    立方正则化 & 梯度与Hessian & Hessian Lipschitz、下有界：$O(\varepsilon^{-3/2})$轮 & 子问题可达$O(d^3)$ \\
    扰动梯度下降 & 梯度与随机扰动 & 高概率返回近似二阶驻点 & 一次梯度，偶尔一次函数值 \\
    \bottomrule
  \end{tabularx}
  \caption{本章主要方法的比较。各行保证的误差对象和条件不同，只能在列出的条件下解读，不构成方法之间的无条件排序。}
  \label{tab:opt-smooth-summary}
\end{table*}

本章始终假设目标可微、没有约束、梯度可以精确地对整个目标求出。后续各章依次放松这些假设：第~\ref{chap:opt-nonsmooth}章处理尖点和由几项组成的复合目标，那里下降引理与邻近映射一起构成新的局部模型；第~\ref{chap:constrained-optimization}章处理可行范围，并在内点法中再次用到牛顿法；第~\ref{chap:opt-stochastic-block}章处理只能抽样得到的梯度和分块更新。多解问题中由数据和算法共同决定选出哪个解，集中在第~\ref{chap:opt-function-fitting}章讨论。
